\chapter{Arc Fault Data Collection Logs}
\label{app:f}

This appendix summarizes the arc-fault data-preparation logs, training report,
and exported embedded-model settings used by the TinyML Smart Plug. The raw
CSV and JSON artifacts remain in the TinyML project directory; this appendix
summarizes the values needed to connect the source artifacts to this paper
without reproducing long console dumps or oversized tables.

\section{Dataset Preparation Summary}
\label{app:f-dataset-preparation}

Table~\ref{tab:app-f-dataset-cleaning} summarizes the cleaned dataset reported
by the data-preparation workflow. The values were taken from the generated
preparation report for the selected training generation.
\begin{table}[H]
  \centering
  \caption{Arc-fault dataset preparation summary}
  \label{tab:app-f-dataset-cleaning}
  \begin{tabular}{lr}
    \toprule
    \textbf{Item} & \textbf{Count} \\
    \midrule
    Raw rows before cleaning & 55,085 \\
    Rows after cleaning & 49,296 \\
    Rows removed & 5,789 \\
    Trainable rows & 45,010 \\
    Trusted normal rows & 17,762 \\
    Conflict-marked rows & 1 \\
    Context-only placeholder rows preserved & 16,353 \\
    Arc-truth available rows & 32,943 \\
    \bottomrule
  \end{tabular}
\end{table}

The cleaning stage removed invalid or unusable rows while preserving rows that
were still needed for context modeling. The resulting set was strongly
imbalanced, which is why this paper reports precision, recall, F1-score,
and event-level metrics instead of relying only on accuracy.

The retained rows were grouped by division and label as shown in
Table~\ref{tab:app-f-label-division}. These counts connect the laboratory
collection procedure in Appendix~\ref{app:e} to the training and evaluation
workflow described in Chapter~\ref{chap:methodology}.
\begin{table}[H]
  \centering
  \caption{Prepared data division and label counts}
  \label{tab:app-f-label-division}
  \begin{tabular}{lr}
    \toprule
    \textbf{Category} & \textbf{Count} \\
    \midrule
    Start-division rows & 18,475 \\
    Steady-division rows & 17,762 \\
    Arc-division rows & 13,059 \\
    Normal-label rows & 49,013 \\
    Arc-label rows & 283 \\
    Arc-training rows & 30,821 \\
    Load-context rows & 18,475 \\
    Conflict policy: none & 49,295 \\
    Conflict policy: prefer more trusted normal & 1 \\
    \bottomrule
  \end{tabular}
\end{table}

The division counts show that the context model was trained from start
segments, while the arc model used the steady and arc divisions. The small
number of arc-labeled rows confirms that series-arc detection was a sparse
event-classification task.

\section{TinyML Training Report and Export Summary}
\label{app:f-tinyml-training-summary}

Table~\ref{tab:app-f-rf-search} summarizes the selected Random Forest settings
from the training report. These settings describe the offline model selected
before export to the embedded firmware.
\begin{table}[H]
  \centering
  \caption{Selected Random Forest training settings}
  \label{tab:app-f-rf-search}
  \begin{tabular}{ll}
    \toprule
    \textbf{Parameter} & \textbf{Selected value} \\
    \midrule
    Model family & Random Forest \\
    Number of trees & 160 \\
    Maximum tree depth & 16 \\
    Minimum samples per leaf & 4 \\
    Minimum samples per split & 12 \\
    Maximum feature fraction & 0.35 \\
    Bootstrap sampling & Enabled \\
    Split criterion & Entropy \\
    Cost-complexity pruning alpha & 0.0001 \\
    Minimum impurity decrease & 0.000001 \\
    Cross-validation splits & 5 \\
    \bottomrule
  \end{tabular}
\end{table}

The final training view contained 35,512 training rows, including 344 positive
rows and 35,168 negative rows. The held-out validation split contained
9,434 rows with 90 positive rows, and the held-out test split contained
10,268 rows with 90 positive rows.

Table~\ref{tab:app-f-rf-test-metrics} gives the held-out test metrics reported
for the offline Random Forest model. These frame-level metrics are the source
of the TinyML model-performance discussion in Chapter~\ref{chap:results}.
\begin{table}[H]
  \centering
  \caption{Held-out frame-level Random Forest test metrics}
  \label{tab:app-f-rf-test-metrics}
  \begin{tabular}{lr}
    \toprule
    \textbf{Metric} & \textbf{Value} \\
    \midrule
    Offline report threshold & 0.25 \\
    Accuracy & 0.997078 \\
    Precision & 0.928571 \\
    Recall & 0.722222 \\
    F1-score & 0.812500 \\
    True negatives & 10,173 \\
    False positives & 5 \\
    False negatives & 25 \\
    True positives & 65 \\
    \bottomrule
  \end{tabular}
\end{table}

Table~\ref{tab:app-f-rf-test-metrics} shows that the frame-level classifier had
high precision but lower recall. This means that false positive frames were
rare, while some positive arc frames were not detected at the frame level.
For this reason, Chapter~\ref{chap:results} also interprets event-level and
trial-level behavior.

The event-level report grouped adjacent positive frames into arc episodes.
Table~\ref{tab:app-f-event-metrics} summarizes this view.
\begin{table}[H]
  \centering
  \caption{Held-out event-level Random Forest test metrics}
  \label{tab:app-f-event-metrics}
  \begin{tabular}{lr}
    \toprule
    \textbf{Metric} & \textbf{Value} \\
    \midrule
    True arc events & 25 \\
    Detected arc events & 21 \\
    Missed arc events & 4 \\
    Predicted events & 31 \\
    Overlapping predicted events & 30 \\
    False-alarm events & 1 \\
    False-alarm sessions & 1 \\
    False alarms per session & 0.1 \\
    Event recall & 0.840000 \\
    Event precision & 0.967742 \\
    Event F1-score & 0.899358 \\
    \bottomrule
  \end{tabular}
\end{table}

The event-level metrics better match the protection problem because a series
arc appears as a short disturbance episode rather than a single independent
frame. The one false-alarm event and four missed events explain why the
prototype should be interpreted as complementary protection and not as a
certified replacement for standard protective devices.

\section{Embedded Model Export Summary}
\label{app:f-embedded-export-summary}

The exported embedded arc model used a context-aware input vector. The source
header defines feature-version 6, an input dimension of 13, six base waveform
features, and seven context fields. The embedded input order is shown in
Table~\ref{tab:app-f-embedded-input-order}.
\begin{table}[H]
  \centering
  \caption{Embedded Random Forest arc-model input order}
  \label{tab:app-f-embedded-input-order}
  \begin{tabular}{rl}
    \toprule
    \textbf{Index} & \textbf{Feature} \\
    \midrule
    0 & Maximum low-current run duration \\
    1 & Zero-dwell ratio \\
    2 & Low-current ratio \\
    3 & Absolute RMS-current z-score versus baseline \\
    4 & Midband residual ratio \\
    5 & Mid/high-frequency spectral flux \\
    6 & Context family: resistive linear \\
    7 & Context family: inductive motor \\
    8 & Context family: rectifier SMPS \\
    9 & Context family: phase-angle controlled \\
    10 & Context family: brush universal motor \\
    11 & Context family: other mixed \\
    12 & Context-family confidence \\
    \bottomrule
  \end{tabular}
\end{table}

The deployed header used an embedded arc threshold of 0.4800 for supported
known-context decisions, a context-confidence minimum of 0.4500, an
unknown-context threshold of 0.7800, and a minimum of three corroborating
arc-like feature votes under unknown context. These embedded settings are
more conservative than the offline report threshold because the deployed
device must limit nuisance relay trips during uncertain load-family context.

The exported context model is a prototype classifier with six inputs:
$\Delta I_{\mathrm{rms}}$, half-cycle asymmetry, midband residual ratio,
current THD, mid/high-frequency spectral flux, and high-frequency energy
delta. It standardizes each input, computes distance to six family centroids,
and accepts a family only when the best context confidence is at least 0.4500.
Otherwise, the arc model is evaluated under the stricter unknown-context
threshold.
